\section{Local cuspidality at three beyond the cubic tower}
\label{sec:localthree}

Two elementary maps to the same proper quotient give local information
that does not require an injection of geometric pro-$3$ groups.
Let $m\ge5$ be odd with $3\nmid m$, let
\[
 C_m:\ v^2+v=u^m,\qquad A_m=\operatorname{Jac}(C_m),
\]
where $C_m$ is the smooth proper model, and embed $C_m$ into $A_m$
using its unique point $\infty$ at infinity.
For a local field $L$, write $A_m(L)^{\wedge3}=\varprojlim_j
A_m(L)/3^jA_m(L)$.

\begin{lemma}[Four points in the zero Kummer locus]\label{lem:fourzero}
The points $Q\in C_m(\mathbb Q_3)$ whose classes $[Q-\infty]$ vanish
in $A_m(\mathbb Q_3)^{\wedge3}$ are exactly
\[
 \infty,\quad (0,0),\quad(0,-1),\quad
 W_m=(a_m,-1/2),\qquad a_m^m=-1/4.
\]
Here $a_m$ is the unique such element of $\mathbb Q_3$.
\end{lemma}
\begin{proof}
The hyperelliptic equation $(2v+1)^2=1+4u^m$ has good reduction at
$3$: its right-hand side is squarefree modulo $3$, and its odd degree
provides the smooth point at infinity. Since $m$ is odd, the map
$u\mapsto u^m$ is the identity on $\mathbb F_3$. Counting one $u$ for
each of the three values of $v$ gives
$\#C_m(\mathbb F_3)=4$, including infinity.

The compact $3$-adic Lie group $A_m(\mathbb Q_3)$ is an extension of
a finite group by an open subgroup isomorphic to a finite direct sum
of $\mathbb Z_3$. Its map to its pro-$3$ completion therefore has as
kernel precisely its finite prime-to-$3$ torsion subgroup. Good
reduction is injective on that subgroup, because the formal kernel
is pro-$3$ and multiplication by an integer prime to $3$ is invertible
there. The Abel--Jacobi map of the special fiber is injective. Thus at
most one zero-Kummer point can lie in any one residue disc, and there
are at most four in total.

The divisors of $v$ and $v+1$ are respectively
$m(0,0)-m\infty$ and $m(0,-1)-m\infty$. These two cusp classes have
order dividing $m$ and hence lie in the kernel. Raising a $3$-adic unit
to the odd power $m$, with $3\nmid m$, is an automorphism: it is an
automorphism on $\{\pm1\}$ and on $1+3\mathbb Z_3$ by the $3$-adic
logarithm. Hence $a_m$ exists uniquely. The point $W_m$ is a branch
point of the hyperelliptic map, so
$\operatorname{div}(u-a_m)=2W_m-2\infty$. Its class also lies in the
kernel. The four displayed points reduce respectively to infinity,
$(0,0)$, $(0,-1)$ and $(-1,1)$, which are distinct. They exhaust the
possible zero-Kummer points.
\end{proof}

\begin{proposition}[A vanishing proper Selmer group forces a local cusp]\label{prop:localtwomaps}
Suppose the $3$-adic Selmer group of $T_3A_m$ is zero; equivalently,
$A_m(\mathbb Q)^{\wedge3}=0$ and $T_3\Sha(\mathbb Q,A_m)=0$.
For every odd multiple $n$ of $m$, the localization at $3$ of any
point-or-cusp Selmer section of the geometrically pro-$3$ abelian
quotient of $U_n/\mathbb Q$ is cuspidal at a rational coordinate
cusp. The same conclusion holds for any geometric pro-$3$ quotient
retaining that abelianization, including the full pro-$3$ group.
\end{proposition}
\begin{proof}
Compose the coordinate-power morphism $X_n\to X_m$ with the two maps
\[
 f_Z(X:Y:Z)=\left(-XY/Z^2,X^m/Z^m\right),\qquad
 f_X(X:Y:Z)=\left(-YZ/X^2,Y^m/X^m\right)
\]
to $C_m$, and then with its Abel--Jacobi embedding. These extend to
morphisms of proper smooth curves. Rational coordinate cusps map to
$\infty,(0,0),(0,-1)$; their classes have order dividing $m$, so their
$3$-adic Kummer classes vanish. Thus projecting the global section
through either morphism gives an ordinary $3$-adic Selmer class of
$T_3A_m$ with respect to the indicated origin. Both classes are zero
by hypothesis. This uses only the geometric abelian quotient of the
global section and its full local point-or-cusp realizations.

If a local realization at $3$ had center in $U_n(\mathbb Q_3)$, its
image $(X:Y:Z)\in U_m(\mathbb Q_3)$ would have all coordinates nonzero.
The preceding lemma applies to both its images on $C_m$. None can be
$\infty,(0,0),(0,-1)$ because the inverse images of these points are
coordinate divisors. Consequently both images must be $W_m$.
The first equation says $X^m/Z^m=-1/2$. The Fermat equation gives
$Y^m/Z^m=-1/2$ as well, hence $Y^m/X^m=1$. The second image would
instead require $Y^m/X^m=-1/2$, a contradiction.

Therefore every local realization at $3$ is a coordinate cusp.
The only roots of unity in $\mathbb Q_3$ are $\{\pm1\}$, so for odd
$n$ the only $n$th root of unity there is $1$. The three local
coordinate cusps are exactly the three rational coordinate cusps.
This proves the stated local conclusion. No conclusion about global
conjugacy follows merely from this argument.
\end{proof}

\begin{proof}[Proof of Theorem~\ref{thm:localthreeintro}]
If $n$ is divisible by $m\in\{5,7,11\}$, the arithmetic
Propositions~\ref{prop:arithmetic}, \ref{prop:septic}, and
\ref{prop:eleven}, respectively, make the proper $3$-adic Selmer
group zero: the rational Mordell--Weil group has order $m$, and the
Tate--Shafarevich group is finite. Apply Proposition~\ref{prop:localtwomaps}.
If $3\mid n$, the coordinate-power map $X_n\to X_3=E$
sends the compact $3$-adic class of any Selmer section to one of
$E(\mathbb Q)=\{O,(3,4),(3,-5)\}$, by Lemma~\ref{lem:cubiclocal}.
The local Kummer map at $3$ is injective, so every local center maps
to that rational cusp. The inverse image of a coordinate cusp under
the power map consists entirely of coordinate cusps. Since
$\mathbb Q_3$ has no nontrivial odd-order roots of unity, the local
center is again one of the three rational coordinate cusps upstairs.
These projections factor through geometric abelianization, proving
the assertion for every quotient in the statement.
\end{proof}


\subsection{Local separation of rational cusp packets}
\label{subsec:lm-affine-packet-separation}

\begin{proposition}[The affine abelianization separates the rational cusps]
\label{prop:lm-affine-ab-cusp-separation}
For every odd $n\ge3$, the three rational cusp packets of
$U_n/\mathbb Q_3$ have pairwise disjoint images in the
geometrically pro-$3$ abelian quotient.
Consequently the three global rational cusp packets are
also pairwise disjoint in that quotient.
If $3\nmid n$, this assertion uses the affine
abelianization: the compact Jacobian Kummer classes of
the rational cusps are all equal.
\end{proposition}
\begin{proof}
If $3\mid n$, compose with the compact cubic map.
The three rational cubic cusp points have distinct
$3$-primary Kummer classes over $\mathbb Q_3$, because
$E(\mathbb Q_3)$ is pro-$3$ and its Kummer map is
injective.  Thus their inverse-image rational cusp
packets are distinct.

Suppose $3\nmid n$.  Choose a prime $q\ge5$ dividing $n$.
The coordinate-power morphism $U_n\to U_q$ carries
distinct rational coordinate cusps to distinct rational
coordinate cusps.  It is enough to separate the packets
on $U_q$; no injectivity of its map on pro-$3$
fundamental groups is required.

By coordinate permutations it is enough to compare
\[
 P=(0:-1:1),\qquad R=(-1:0:1).
\]
Put $L=\mathbb Q_3(\zeta_q)$, with $\zeta_q$ primitive,
and consider the rational unit on $U_{q,L}$
\[
 \varphi=\frac{Y+\zeta_q Z}{Y+\zeta_q^2 Z}.
\]
Its divisor on the compact Fermat curve is
$q(Q_{\zeta_q}-Q_{\zeta_q^2})$, where
$Q_\zeta=(0:-\zeta:1)$.  In particular $\varphi$
is regular and nonzero at both $P$ and $R$, with
\[
 \varphi(P)=\frac1{1+\zeta_q},
 \qquad
 \varphi(R)=\zeta_q^{-1}.
\]
For any cusp section centered at $P$ or $R$, the
map induced by $\varphi$ to $\mathbb G_m$ is the
Kummer section of the corresponding displayed
value.  The tangent parameter has no effect,
because $\varphi$ has order zero at that cusp.

Equality of two abelian cusp sections at $P$ and
$R$, after restriction to $L$, would therefore give
\[
 \left[\frac{\zeta_q}{1+\zeta_q}\right]=0
 \quad\text{in}\quad H^1(L,\mathbb Z_3(1)).
\]
For a finite extension $L/\mathbb Q_3$, the kernel
of $L^\times\to (L^\times)^{\wedge_3}$ is exactly
the finite subgroup of roots of unity of order
prime to $3$.  The displayed algebraic number is
not a root of unity: choose a complex embedding
taking $\zeta_q$ to $\exp(2\pi i/q)$, under which
its absolute value is
\[
 \frac1{2\cos(\pi/q)}<1
 \qquad(q\ge5).
\]
This is a contradiction.  Coordinate permutations
separate each pair of rational cusps.

Finally, the difference of any two coordinate
cusps on $X_n$ is killed by $n$, using the ratio
of their total-flex tangent lines.  When $3\nmid n$,
these differences have zero $3$-primary Kummer
class.  This verifies the asserted distinction
between compact and affine abelianization.
Global equality would imply local equality at
$3$, so the global packets are disjoint as well.
\end{proof}

\begin{remark}[Scope of local packet separation]\label{rem:affine-packet-scope}
The proposition identifies a unique rational cusp
from a known local abelian cusp class at $3$.
It does not imply that a global abelian Selmer
section is globally cuspidal.  For mixed exponents, the nonzero relative
localization kernel in Proposition~\ref{prop:mixed-kernel} remains,
and membership in this kernel does not itself
give a local geometric representative at the
other places.  Both distinctions are necessary:
compact ambiguity of the cusp center is repaired
by the proposition, whereas the global Selmer
obstruction is not.
\end{remark}

